% TOP_PROBLEM: {"id": 20001085, "title": "Exact quadratic-residue half-set counts and a quantitative Varnavides transfer", "category_id": 2, "category": {"id": 2, "name": "combinatorics", "display_name": "Combinatorics", "description": "Counting problems, graph theory, discrete structures.", "slug": "combinatorics", "order_index": 2, "created_at": "2026-07-31T15:26:25.671Z"}, "status": "partially_solved", "problem_number": "AIM-COMBINATORICS-0210", "set_id": 14, "statusQualification": "Counts ordered nonconstant progressions and rounds the density cardinality down; these conventions resolve counting and integrality ambiguities."}
% FIRST_POSTED: 2026-10-01
% ENTRY_KIND: statement_only
% UnsolvedMath ID 20001085; source catalogue code AIM-COMBINATORICS-0210
% !TeX program = lualatex
% Definition and sources only; this file is not an LLM solution attempt.
\documentclass[11pt,a4paper]{article}
\usepackage[margin=25mm]{geometry}
\usepackage{fontspec}
\setmainfont{lmroman10-regular.otf}[BoldFont=lmroman10-bold.otf,
 ItalicFont=lmroman10-italic.otf,BoldItalicFont=lmroman10-bolditalic.otf]
\usepackage{amsmath,amssymb,mathtools}
\usepackage{unicode-math}
\setmathfont{latinmodern-math.otf}
\AtBeginDocument{\let\setminus\smallsetminus}
\usepackage{xurl,hyperref}
\hypersetup{colorlinks=true,urlcolor=blue}
\setcounter{secnumdepth}{0}
\setlength{\parindent}{0pt}
\setlength{\parskip}{5pt}
\setlength{\emergencystretch}{3em}
\begin{document}
\section*{Exact quadratic-residue half-set counts and a quantitative Varnavides transfer}\label{20001085}
{\small UnsolvedMath ID 20001085 · AIM-COMBINATORICS-0210 · Combinatorics · AIM Workshop Problem Lists}
\subsection{Definitions and mathematical statement}
Let $p$ be an odd prime and $A\subseteq\mathbb F_p$. Count nonconstant
three-term progressions by
\[
 T_p(A)=\bigl|\{(a,d)\in\mathbb F_p\times\mathbb F_p^*:
                              a,a+d,a+2d\in A\}\bigr|.
\]
All arithmetic takes place modulo $p$. The common difference $d$
is nonzero, so all three terms are distinct. This is a count of
ordered progressions: both directions are counted, and in the small
field $\mathbb F_3$ a three-element set has more than one choice of
middle term. Defining the count by $(a,d)$ removes such multiplicity
ambiguities. Constant progressions are excluded.
For an integer $0\le m\le p$, put
\[
                         M(p,m)=\min_{|A|=m}T_p(A).
\]
The first question is to determine $M(p,(p-1)/2)$, or obtain its
sharp asymptotic behaviour as $p\to\infty$ through odd primes.
The second asks the corresponding question for density below one
half: for fixed $0<\delta<1/2$, determine
$M(p,\lfloor\delta p\rfloor)$, including the dependence on $\delta$.
The floor makes the requested cardinality integral without changing
its limiting density.
A useful normalized target is $M(p,\lfloor\delta p\rfloor)/p^2$;
the workshop remarks specifically motivate understanding its small
$\delta$ behaviour. All subsets of the prescribed size are eligible,
without algebraic or symmetry restrictions, so calculations for one
family alone do not identify the minimum.
\subsection{Short English statement}
Among subsets of a prime field of a prescribed density at or below one half, how few nonconstant three-term progressions can occur?
\subsection{Sources}
\begin{itemize}
\item[S1] ULAM.AI and the UnsolvedMath Contributors, supplied problems.json, record 20001085 (AIM-COMBINATORICS-0210), AIM Workshop Problem Lists. Catalogue identity and source transcription; CC BY 4.0.
\url{https://huggingface.co/datasets/ulamai/UnsolvedMath}.
\item[S2] American Institute of Mathematics, Recent trends in additive combinatorics, item 1.4. Original workshop question underlying AIM-COMBINATORICS-0210.
\url{https://aimath.org/WWN/additivecomb/additivecomb.pdf}.
\end{itemize}
\end{document}
